\documentclass[11pt]{scrartcl}

\usepackage{amsthm}
\usepackage{hyperref}
\usepackage{lastpage}
\usepackage{minted}

\theoremstyle{definition}
\newtheorem*{example}{Example} 

\usepackage[none]{hyphenat}
\setlength{\parindent}{0pt}
% setlength{\emergencystretch}{3em}

\renewcommand\labelitemi{$\vcenter{\hbox{\tiny$\bullet$}}$} 

\usepackage[margin=1in]{geometry}
\usepackage{fancyhdr}
\pagestyle{fancy}
\lhead{ESC180: Introduction to Computer Programming} 
\rhead{\thepage}
\cfoot{}

\title{ESC180: Introduction to Computer Programming \\[0.1in] Exam Review}
\author{Brendan Tang}
\date{Fall 2026}

\begin{document}
\maketitle
\thispagestyle{empty}

\tableofcontents

\newpage
\section{Python Fundamentals}
\subsection{Comments and Output}

We can use the \mintinline{python}{print()} function to display text or output values. Text must be inside quotes, 
otherwise Python will return a syntax error.

\begin{example}
    \
\begin{minted}{python} 
    print("Hello world!") 
    print("I love Praxis!") # Automatically prints on a new line
\end{minted} 
    \textit{Output.}
\begin{minted}{text}
    Hello world!
    I love Praxis!
\end{minted}
\end{example}

Any text after \mintinline{python}{#} is a comment. These will not be read by Python, and are used to make code
more readable. You can also do math inside of a print statement. Observe the absence of quotes.

\begin{example}
    \
\begin{minted}{python}
    print("What is 3 + 10?") 
    print("3 + 10 =", 3 + 10) # Prints 13, as 3 + 10 = 13
\end{minted}
    \textit{Output.}
\begin{minted}{text}
    What is 3 + 10?
    3 + 10 = 13
\end{minted}
\end{example}

\subsection{Variables and Data Types}

Here are some commonly used data types in Python.

\begin{center}
\begin{tabular}{|l|l|} \hline
    \textbf{Data Type} & \textbf{Example} \\ \hline
    str & \mintinline{python}{"hi!"} \\ \hline
    int & \mintinline{python}{120} \\ \hline
    float & \mintinline{python}{12.5} \\ \hline
    bool & \mintinline{python}{True} \\ \hline
    list & \mintinline{python}{["an", "ban", "can"]} \\ \hline
    tuple & \mintinline{python}{("an", "ban", "can")} \\ \hline
    dict & \mintinline{python}{{"ESC101": 70, "ESC194": 100}} \\ \hline
\end{tabular}
\end{center}

Variables are used to store data values. A variable is created when you assign a value to it with the assignment
operator, \mintinline{python}{=}.

\begin{example}
    \
\begin{minted}{python}
    toggled = False
    name = "John Doe"

    x1 = 50 # stores 50
    x2 = float(50) # stores 50.0

    y1 = 11.1 # stores 11.1
    y2 = int(11.1) # stores 11
\end{minted}
\end{example}

Observe how Python automatically determines the type of a variable, unless you explicitly cast a different data
type.

\subsubsection{Interlude (Local and Global Variables)}

Variables that are created outside of a function are known as global variables. Global variables can be used by 
everyone, both inside of functions and outside. Variables that are created inside a function are known as local 
variables. Local variables can only be used inside their functions. 

We can however create global variables inside a function by using the \mintinline{python}{global} keyword.

\subsection{Booleans and Operators}

Comparing two values generates a Boolean value, \mintinline{python}{True} or \mintinline{python}{False}. You can
also use \mintinline{python}{bool()}, which will return \mintinline{python}{True} unless you pass in numeric zeros, 
empty sequences/collections, or \mintinline{python}{None} in which case \mintinline{python}{False} is returned.

Here are some arithmetic operators.

\begin{center}
\begin{tabular}{|l|l|} \hline
    \textbf{Arithmetic Operator} & \textbf{Name} \\ \hline
    \mintinline{python}{+} & Addition \\ \hline
    \mintinline{python}{-} & Subtraction \\ \hline
    \mintinline{python}{*} & Multiplication \\ \hline
    \mintinline{python}{/} & Division \\ \hline
    \mintinline{python}{%} & Modulus \\ \hline
    \mintinline{python}{**} & Exponentiation \\ \hline
    \mintinline{python}{//} & Floor division \\ \hline
\end{tabular}
\end{center}

Here are some assignment operators.

\begin{center}
\begin{tabular}{|l|l|l|} \hline
    \textbf{Assignment Operator} & \textbf{Example} & \textbf{Equivalent To} \\ \hline
    \mintinline{python}{=} & \mintinline{python}{x = 10} & \mintinline{python}{x = 10} \\ \hline
    \mintinline{python}{+=} & \mintinline{python}{x += 10} & \mintinline{python}{x = x + 10} \\ \hline
    \mintinline{python}{-=} & \mintinline{python}{x -= 10} & \mintinline{python}{x = x - 10} \\ \hline
    \mintinline{python}{*=} & \mintinline{python}{x *= 10} & \mintinline{python}{x = x * 10} \\ \hline
    \mintinline{python}{/=} & \mintinline{python}{x /= 10} & \mintinline{python}{x = x / 10} \\ \hline
    \mintinline{python}{%=} & \mintinline{python}{x %= 10} & \mintinline{python}{x = x % 10} \\ \hline
    \mintinline{python}{**=} & \mintinline{python}{x **= 10} & \mintinline{python}{x = x ** 10} \\ \hline
    \mintinline{python}{//=} & \mintinline{python}{x //= 10} & \mintinline{python}{x = x // 10} \\ \hline
\end{tabular}
\end{center}

Here are some comparison operators.

\begin{center}
\begin{tabular}{|l|l|} \hline
    \textbf{Comparison Operator} & \textbf{Name} \\ \hline
    \mintinline{python}{==} & Equal \\ \hline
    \mintinline{python}{!=} & Not equal \\ \hline
    \mintinline{python}{>} & Greater than \\ \hline
    \mintinline{python}{<} & Less than \\ \hline
    \mintinline{python}{>=} & Greater than or equal to \\ \hline
    \mintinline{python}{<=} & Less than or equal to \\ \hline
\end{tabular}
\end{center}.

Here are some logical, identity, and membership operators. \vspace{-0.5\baselineskip} % Help pls
\begin{center}
\begin{tabular}{|l|l|} \hline
    \textbf{Operator} & \textbf{Description} \\ \hline
    \mintinline{python}{and} & Returns \mintinline{python}{True} if both statements are 
    \mintinline{python}{True} \\ \hline
    \mintinline{python}{or} & Returns \mintinline{python}{True} if one of the statements is 
    \mintinline{python}{True} \\ \hline
    \mintinline{python}{not} & Reverses the result, returns \mintinline{python}{False} if the result is 
    \mintinline{python}{True} \\ \hline
    \mintinline{python}{is} & Returns \mintinline{python}{True} if both variables are the same object \\ \hline
    \mintinline{python}{is not} & Returns \mintinline{python}{True} if both variables are not the same 
    object \\ \hline
    \mintinline{python}{in} & Returns \mintinline{python}{True} if a sequence with the specified value is present 
    in the object \\ \hline
    \mintinline{python}{not in} & Returns \mintinline{python}{True} if a sequence with the specified value is not 
    present in the object \\ \hline
\end{tabular}
\end{center}

\subsection{Conditional Statements}

The \mintinline{python}{if} statement evaluates an expression to determine its boolean value. If 
\mintinline{python}{True}, the code block inside the statement is executed. If \mintinline{python}{False}, the 
code block is skipped.

\begin{example}
    \
\begin{minted}{python}
    if grade >= 85:
        print("GPV: 4.0")
    elif grade >= 80: # Python will move on to this if the original is False
        print("GPV: 3.7")
    elif grade >= 77: # Same for this, if both of the above are False
        print("GPV: 3.3")
    else: # Code block will be executed if all of the above are False
        print("Is this Praxis?")
\end{minted}
\end{example}

Notice that in Python, indentation is used to determine which statements belong together. \\

If we have a lot of cases, we may wish to use a \mintinline{python}{match} statement instead.

\begin{example}
    \
\begin{minted}{python}
    match color:
        case "red":
            print("Roses are red.")
        case "yellow" | "green": # We can check for more than one value match
            print("Grass is green, yellow if dead.")
        case "blue": if time == "day" # We can add an extra condition check
            print("The sky is blue.")
        case _: # This will execute if there are no matches
            print("This color is boring.")
\end{minted}
\end{example}

\subsection{Loops}

A \mintinline{python}{for} loop is used for iterating over a sequence (a list, a tuple, a set, or a string).
Most commonly, we use the \mintinline{python}{range} data type, which represents an immutable sequence of 
numbers.

\begin{example}
    \
\begin{minted}{python}
    for i in range(3): # Notice how by default, we begin at 0
        print("I am", i)

    print(" ")

    for i in range(3, 10, 3): # Syntax is (start, stop, step)
        print("I am", i)
\end{minted}
    \textit{Output.}
\begin{minted}{text}
    I am 0
    I am 1
    I am 2
     
    I am 3
    I am 6
    I am 9
\end{minted}
\end{example}

A \mintinline{python}{while} loop will execute a set of statements as long as a condition is 
\mintinline{python}{True}. If the condition is always \mintinline{python}{True}, we could end up in an 
infinite loop. 

\begin{example}
    \
\begin{minted}{python}
    height = 1
    while True: # This loop will continue forever, unless...
        print("I am", height, "feet tall!")
        height += 1
        if height > 6:
            break # This will exit the loop
\end{minted}
    \textit{Output.}
\begin{minted}{text}
    I am 1 feet tall!
    I am 2 feet tall!
    I am 3 feet tall!
    I am 4 feet tall!
    I am 5 feet tall!
    I am 6 feet tall!
\end{minted}
\end{example}

Observe the use of \mintinline{python}{break} above. If we wish to stop the current iteration but stay in the 
loop, we can use \mintinline{python}{continue}. We can also use an \mintinline{python}{else} statement to run 
a block of code when the condition is no longer \mintinline{python}{True}. This means that exiting the loop 
using \mintinline{python}{break} will not run the code inside the \mintinline{python}{else} statement.

\subsection{Strings}

Strings in Python are arrays of unicode characters. However, since Python doesn't have a character data type, a 
character is simply a string with length 1. \\

Here are some common string methods. Note that there are many others, these are just the most common. 

\begin{center}
\begin{tabular}{|l|l|} \hline
    \textbf{Syntax} & \textbf{Description} \\ \hline
    \mintinline{python}{s[i]} & Get the character at position \mintinline{python}{i} \\ \hline
    \mintinline{python}{len(s)} & Returns the length of \mintinline{python}{s} \\ \hline
    \mintinline{python}{str0 in s} & Checks if \mintinline{python}{str0} is in \mintinline{python}{s} 
    \\ \hline
    \mintinline{python}{s[start:stop:step]} & Extracts a portion of \mintinline{python}{s} \\ \hline
    \mintinline{python}{s.upper()} & Returns \mintinline{python}{s} in upper case \\ \hline
    \mintinline{python}{s.lower()} & Returns \mintinline{python}{s} in lower case \\ \hline
    \mintinline{python}{s.strip()} & Returns \mintinline{python}{s} without leading/trailing whitespace \\ \hline
    \mintinline{python}{s.replace(str1, str2)} & Returns \mintinline{python}{s}, replacing 
    \mintinline{python}{str1} with \mintinline{python}{str2} \\ \hline
    \mintinline{python}{s.split(char)} & Returns a list of substrings, with separator \mintinline{python}{char} 
    \\ \hline
\end{tabular}
\end{center}

Python allows negative indexing. Also, since strings are arrays, we can loop through the characters in a string 
with a \mintinline{python}{for} loop.

\begin{example}
    \
\begin{minted}{python}
    i = 1
    for char in "mathematics":
        print((char * i) + (" " * 2 * (11 - i)) + (char * i))
        i += 1
\end{minted}
\newpage % dang it dang it dang it :( this makes me very sad
    \textit{Output.}
\begin{minted}{text} 
    m                    m
    aa                  aa
    ttt                ttt
    hhhh              hhhh
    eeeee            eeeee
    mmmmmm          mmmmmm
    aaaaaaa        aaaaaaa
    tttttttt      tttttttt
    iiiiiiiii    iiiiiiiii
    cccccccccc  cccccccccc
    ssssssssssssssssssssss
\end{minted} 
\end{example}

\subsection{Functions}

A function is defined using the \mintinline{python}{def} keyword. To call a function, write its name followed by 
parentheses.

\begin{example}
    \
\begin{minted}{python}
    def print_hello():
        print("Hello there!")
    
    print_hello()
\end{minted}
    \textit{Output.}
\begin{minted}{text}
    Hello there!
\end{minted}
\end{example}

Functions can send data back to the code that called them using the \mintinline{python}{return} statement. When 
a function reaches a \mintinline{python}{return} statement, it stops executing and sends the result back. If a 
function doesn't have a \mintinline{python}{return} statement, it returns \mintinline{python}{None} by default. \\

We can also pass information into functions as arguments. These are specified after the function name, inside the 
parentheses. You can also assign default values to parameters.

\begin{example}
    \
\begin{minted}{python}
    def greeting_message(name1, name2 = "alien"): # name2 has a default value
        return "Hi " + name1 + ", hello " + name2 + "!"
    
    greeting_message("cat", "dog") # Nothing is printed (why?)
    print(greeting_message("cow")) # Here we exclude a second argument
\end{minted}
    \textit{Output.}
\begin{minted}{text}
    Hi cow, hello alien!
\end{minted}
\end{example}

If you do not know how many arguments will be passed into your function, add a \mintinline{python}{*} before the 
parameter name to send a tuple of arguments. You can also add a \mintinline{python}{**} before the parameter name 
to send the function a dictionary of arguments.

\subsection{User Input}
You can use \mintinline{python}{input()} to ask the user for input. You can include a prompt inside the parentheses, which 
is a message that will display before the user input on the same line. The input from the user is always 
treated as a string.

\subsection{Lists and Tuples}

Lists are used to store multiple items in a single variable. List items are indexed, the first item has 
\mintinline{python}{[0]}, the second item has index \mintinline{python}{[1]}, etc. A list can contain different 
data types. 

Here are some common list methods.

\begin{center}
\begin{tabular}{|l|l|} \hline
    \textbf{Syntax} & \textbf{Description} \\ \hline
    \mintinline{python}{l[i]} & Get the item at position \mintinline{python}{i} in \mintinline{python}{l} \\ \hline
    \mintinline{python}{len(l)} & Returns the length of \mintinline{python}{l} \\ \hline
    \mintinline{python}{l[start:stop:step]} & Extracts a portion of \mintinline{python}{l} \\ \hline
    \mintinline{python}{l.insert(i, item)} & Inserts \mintinline{python}{item} in \mintinline{python}{l} 
    at \mintinline{python}{i} \\ \hline
    \mintinline{python}{l.append(item)} & Appends \mintinline{python}{item} to \mintinline{python}{l} \\ \hline
    \mintinline{python}{list1.extend(list2)} & Appends \mintinline{python}{list1} to \mintinline{python}{list2} 
    \\ \hline
    \mintinline{python}{l.remove(item)} & Removes the first occurence of \mintinline{python}{item} from 
    \mintinline{python}{l} \\ \hline
    \mintinline{python}{l.pop(i)} & Removes \mintinline{python}{i} from \mintinline{python}{l} \\ \hline
    \mintinline{python}{l.clear()} & Empties \mintinline{python}{l} \\ \hline
    \mintinline{python}{l.sort()} & Sorts \mintinline{python}{l} in ascending order \\ \hline
    \mintinline{python}{l.sort(reverse = True)} & Sorts \mintinline{python}{l} in descending order \\ \hline
    \mintinline{python}{l.reverse()} & Reverses the order of \mintinline{python}{l} \\ \hline
    \mintinline{python}{l.copy()} & Copies \mintinline{python}{l} \\ \hline
\end{tabular}
\end{center}

\subsection{Dictionaries}

Dictionaries are used to score data values in key:value pairs. I don't feel like making another table so here is 
an example.

\begin{example}
    \
\begin{minted}{python} 
    grades = {
        "Praxis": 60,
        "Calculus": 100,
        "Civ": 80
    }
    print("My grade in Praxis is " + grades.get("Praxis"))

    if "Civ" in grades: # Checks if "Civ" is a key in grades
        print("I am taking CIV102 this term!") 
    else:
        print("I am not taking CIV102 this term.")

    grades.pop("Civ") # Removes the pair with "Civ" 

    if "Civ" in grades: # Checks if "Civ" is a key in grades
        print("I am taking CIV102 this term!")
    else:
        print("I am not taking CIV102 this term.")

    grades.clear() # Empties the dictionary
\end{minted} 
    \textit{Output.}
\begin{minted}{text}
    My grade in Praxis is 60
    I am taking CIV102 this term!
    I am not taking CIV102 this term.
\end{minted}
\end{example}

\end{document}